\documentclass[10pt,a4paper]{amsart}
\usepackage[margin=19mm]{geometry}
\usepackage{amsmath,amssymb,mathtools}
\usepackage[scr=rsfs,cal=euler]{mathalfa}
\usepackage{xfrac}
\usepackage[unicode,hidelinks,bookmarks=false]{hyperref}
\newcommand{\ie}{i.e.\ }
\newcommand{\cf}{cf.\ }
\newcommand{\resp}{respectively\ }
\newcommand{\Subsection}{\subsection}
\providecommand{\nobreakdash}{\nobreak\hbox{-}}
\setlength{\emergencystretch}{2em}
\multlinegap0pt
\usepackage{bm}
\usepackage{bbm}
\usepackage{dsfont}
\usepackage{xspace}
\usepackage{cancel}
\usepackage{mathtools}
\usepackage{amsthm}
\makeatletter
\@ifundefined{theorem}{\newtheorem{theorem}{Theorem}}{}
\@ifundefined{lemma}{\newtheorem{lemma}[theorem]{Lemma}}{}
\@ifundefined{proposition}{\newtheorem{proposition}[theorem]{Proposition}}{}
\makeatother
\newcommand{\Ncal}{\mathcal N}
\newcommand{\dd}{\mathrm d}
\title[Nilpotent Jacobian maps in dimension three and stable tamene]{Nilpotent Jacobian maps in dimension three and stable tameness in block extensions}
\author{Jie Wang, Dan Yan}
\date{Source: arXiv:2609.16580v1 (CC BY 4.0)}
\begin{document}
\maketitle

\noindent\emph{Source and licence.} Nilpotent Jacobian maps in dimension three and stable tameness in block extensions. Version arXiv:2609.16580v1 (CC BY 4.0), distributed under the Creative Commons Attribution 4.0 International licence, as recorded in the arXiv licence metadata for this exact version and shown on the abstract page https://arxiv.org/abs/2609.16580v1. Full rights notice: licenses/arxiv-2609.16580-CC-BY-4.0.txt.\par\smallskip

\noindent\textit{Editorial scope.} Counted unit: the Proposition whose source label is \texttt{prop:pair-normal} in the article (source0.tex lines 415--417, source label \texttt{prop:pair-normal}), delivered together with the article's own complete original proof of it (source0.tex lines 419--569, one \texttt{proof} environment of 151 lines). No mathematics is written, repaired or re-derived: statement and proof are the article's own text line for line. Required context retained uncounted: lem:field-facts (lines 186--193). Every definition, display and label of those lines is retained unchanged. Omitted from this package and not redistributed: 1--185, 194--414, 418--418, 570--936. Nothing inside a retained excerpt is omitted or altered: every retained body is delivered exactly as the article's lines are written, so no omission receipt of any kind accompanies this package. Citations: the retained excerpts carry no citation command at all, so no bibliography entry of the article is transcribed and no reference is redistributed. Reference adaptation: none was needed and none was made; every reference inside the delivered unit resolves inside it, so no unresolved reference or citation is produced. Printed slips kept literal: no printed word, symbol, hypothesis or number of the counted statement or of its proof was altered. Layout and class: the article's own class is \texttt{amsart} with packages including fontenc, lmodern, amsmath, amssymb, amsthm, mathtools, microtype, geometry, enumitem, xcolor, hyperref; the wrapper is the lane's pinned \texttt{amsart} editorial wrapper at 10pt on A4, which restates the theorem-like environments the retained text needs and adds the article's own macro definitions that the retained text needs (\texttt{\textbackslash Ncal}, \texttt{\textbackslash dd}), with extra wrapper packages (bm, bbm, dsfont, xspace, cancel, mathtools). The delivered numbering is the unit's own. Assets: no retained span contains a figure environment, a graphics-inclusion command, a PDF-page inclusion, a table or a TikZ picture; no image file is copied into this package, the package declares no asset dependency and no image file is redistributed with it. Licence: version arXiv:2609.16580v1 of this article is distributed under the Creative Commons Attribution 4.0 International licence, as recorded in the arXiv licence metadata for this exact version and shown on the abstract page https://arxiv.org/abs/2609.16580v1. A full rights notice is delivered at licenses/arxiv-2609.16580-CC-BY-4.0.txt. Characters: every retained excerpt is pure ASCII, so the delivered engine, XeLaTeX with its default Latin Modern text font, reports no missing character.
\begin{lemma}\label{lem:field-facts}
The following statements hold in characteristic zero.
\begin{enumerate}
\item If $\mathcal K/E$ is finitely generated and $L$ is the relative algebraic closure of $E$ in $\mathcal K$, then $L/E$ is finite and $\mathcal K/L$ is regular. Every finitely generated $L$-subalgebra of $\mathcal K$ is geometrically integral.
\item A derivation of a field extends uniquely to a finite separable extension. The extensions of commuting derivations commute.
\item Suppose $\mathcal K/k(t_1,\ldots,t_m)$ is finite separable, and extend the coordinate derivations to $\mathcal K$. If $\bar{a}\in \mathcal K$ is killed by $\partial_{t_i}$ for $i\in I$, then $\bar{a}$ is algebraic over $k(t_j:j\notin I)$.
\end{enumerate}
\end{lemma}

\begin{proposition}\label{prop:pair-normal}
Let $H\in K[x,y,z]^3$ be normalized, with nilpotent Jacobian and linearly independent components. If two independent constant linear combinations of its components are algebraically dependent over $K$, then $H$ is linearly conjugate over $K$ to $\Ncal_P$ for some $P\in K[t]\setminus K$ with $P(0)=0$.
\end{proposition}

\begin{proof}
A constant conjugation and the closed-generator theorem give
\begin{equation}\label{eq:pair-form}
 H=(f(r),v,g(r)),\qquad r(0)=f(0)=g(0)=0,
\end{equation}
where $r\in K[x,y,z]$ is closed and $f,g\in K[t]$ are linearly independent. Their derivatives are likewise linearly independent, so
\begin{equation}\label{eq:a-ratio}
 a(t)=\frac{g'(t)}{f'(t)}\in K(t)\setminus K,\qquad a'(t)\neq0.
\end{equation}
Primes in this proof denote univariate derivatives.

Put $W=f'(r)r_x+g'(r)r_z$. The determinant vanishes because the first and third Jacobian rows are proportional. The other nilpotence equations are
\begin{equation}\label{eq:pair-nilpotence}
 W+v_y=0,\qquad Wv_y-r_y\bigl(f'(r)v_x+g'(r)v_z\bigr)=0.
\end{equation}
We claim that $r_y\neq0$. If $r_y=0$, the two equations give $W=v_y=0$. Reducing $W=0$ modulo $r-\lambda$ gives
\[
 r-\lambda\mid f'(\lambda)r_x+g'(\lambda)r_z.
\]
The polynomial on the right has degree less than $\deg r$, and hence is zero. Two choices of $\lambda\in K$ with independent vectors $(f'(\lambda),g'(\lambda))$ force $r_x=r_z=0$, a contradiction.

\smallskip\noindent\emph{A derivation with a slice.}
On $\mathcal K=K(x,y,z)$ define
\begin{equation}\label{eq:delta}
 \Delta=\partial_x+a(r)\partial_z+\frac{v_y}{f'(r)r_y}\partial_y.
\end{equation}
Equations~\eqref{eq:pair-nilpotence} give
\[
 \Delta x=1,\qquad \Delta r=\Delta v=0,\qquad \Delta z=a(r).
\]
Set $s=z-a(r)x$. Then $\Delta s=0$, and $\det J(x,r,s)=r_y\neq0$. Thus $\mathcal K/K(x,r,s)$ is finite separable. Extend $\partial_r,\partial_s$ to $\mathcal K$. They commute with $\Delta$, the extension of $\partial_x$ in these coordinates. If $M=\ker_{\mathcal K}\Delta$, then both preserve $M$.

Evaluating on $x,r,s$ and using uniqueness of extension gives
\begin{equation}\label{eq:dy-identity}
 \partial_y=r_y\bigl(\partial_r-a'(r)x\partial_s\bigr).
\end{equation}
Since $v\in M$,
\[
 \Delta y=\frac{v_r-a'xv_s}{f'},
\]
where $a',f'$ are evaluated at $r$. Subtracting a polynomial in $x$ and applying $\Delta$ gives
\begin{equation}\label{eq:y-polynomial}
 y=c_0+\frac{v_r}{f'}x-\frac{a'v_s}{2f'}x^2,\qquad c_0\in M.
\end{equation}
Together with $z=s+a(r)x$, this proves $K[x,y,z]\subseteq M[x]$. The variable $x$ is transcendental over $M$: differentiating a minimal polynomial would otherwise contradict $\Delta x=1$ and separability.

\smallskip\noindent\emph{A unit calculation.}
Apply~\eqref{eq:dy-identity} to $y$:
\begin{equation}\label{eq:unit-identity}
 1=r_y\bigl(\partial_r-a'x\partial_s\bigr)y.
\end{equation}
Both factors are in $M[x]$. For the second factor, use~\eqref{eq:y-polynomial} and the fact that the extended derivations preserve $M$. Both factors are therefore units of $M[x]$. The coefficient of $x^3$ in the second factor is
\[
 \frac{(a')^2}{2f'}v_{ss},
\]
and therefore $v_{ss}=0$. The elements $v_s$ and $v-sv_s$ are killed by $\Delta$ and $\partial_s$. Lemma~\ref{lem:field-facts} shows that they are algebraic over $K(r)$. Since $r$ is closed, they lie in $K(r)$. Consequently
\begin{equation}\label{eq:v-rational}
 v=\alpha(r)z+\beta(r)x+\gamma(r),\qquad \alpha,\beta,\gamma\in K(t).
\end{equation}
More precisely, $\beta=-a\alpha$, although this relation is not needed in the remaining reduction.

We next show that the coefficients in~\eqref{eq:v-rational} are polynomials. Write
\[
 d(r)v=\alpha_0(r)z+\beta_0(r)x+\gamma_0(r)
\]
with $d,\alpha_0,\beta_0,\gamma_0\in K[t]$ have greatest common divisor one. If $d$ is nonconstant, choose a root $\lambda$ over $\overline K$. Reduction modulo $r-\lambda$ shows that $r-\lambda$ divides
$\alpha_0(\lambda)z+\beta_0(\lambda)x+\gamma_0(\lambda)$. The divisor has positive $y$-degree, whereas the dividend is independent of $y$. The dividend is therefore zero. All four univariate polynomials then share the root $\lambda$, contradicting their B\'ezout identity. Thus
\begin{equation}\label{eq:v-polynomial}
 \alpha,\beta,\gamma\in K[t].
\end{equation}

\smallskip\noindent\emph{Commuting affine derivations.}
Substitute~\eqref{eq:v-rational} into the trace equation. We obtain
\begin{equation}\label{eq:affine-trace}
 f'(r)r_x+g'(r)r_z+
 \bigl(\alpha'(r)z+\beta'(r)x+\gamma'(r)\bigr)r_y=0.
\end{equation}
For $\lambda\in K$, let
\[
 E_\lambda=f'(\lambda)\partial_x+g'(\lambda)\partial_z+
 \bigl(\alpha'(\lambda)z+\beta'(\lambda)x+\gamma'(\lambda)\bigr)\partial_y.
\]
Reduction of~\eqref{eq:affine-trace} modulo $r-\lambda$ gives $r-\lambda\mid E_\lambda r$. Since $\deg E_\lambda r\leq\deg r$, we have $E_\lambda r=\mu_\lambda(r-\lambda)$ for some $\mu_\lambda\in K$. Assign weights $1,2,1$ to $x,y,z$. The derivation $E_\lambda$ strictly lowers weighted degree, so it is locally nilpotent. This forces $\mu_\lambda=0$, so $E_\lambda r=0$ for every $\lambda\in K$.

Choose two values of $\lambda$ for which $(f'(\lambda),g'(\lambda))$ are independent. Constant linear combinations of the corresponding derivations give
\[
 D_1=\partial_x+b_1(x,z)\partial_y,\qquad
 D_2=\partial_z+b_2(x,z)\partial_y,
\]
with $b_1,b_2$ affine and $D_1r=D_2r=0$. Their commutator is
\[
 [D_1,D_2]=(\partial_xb_2-\partial_zb_1)\partial_y.
\]
Its coefficient is constant. Since the commutator annihilates $r$ and $r_y\neq0$, its coefficient is zero. There is therefore a polynomial $Q\in K[x,z]$ with
\[
 \deg Q\leq2,\quad Q(0)=0,\qquad Q_x=-b_1,\quad Q_z=-b_2.
\]
In the auxiliary polynomial coordinates $(x,\rho,z)$, where $\rho=y+Q(x,z)$, the derivations become $\partial_x,\partial_z$. Their common kernel in the polynomial ring is $K[\rho]$, so $r=R(\rho)$ for a polynomial $R$ with $R(0)=0$. Absorb $R$ into the univariate polynomials in~\eqref{eq:pair-form} and~\eqref{eq:v-rational}.

Write
\[
 Q=\tfrac12Ax^2+Bxz+\tfrac12Cz^2+Dx+Ez.
\]
Expressing the trace equation in the coordinates $x,\rho,z$ and comparing coefficients gives
\[
 \beta'=-Af'-Bg',\qquad \alpha'=-Bf'-Cg',\qquad \gamma'=-Df'-Eg'.
\]
Integrating and using $H(0)=0$, we obtain
\begin{equation}\label{eq:quadratic-form}
 H=\bigl(f(\rho),-Q_xf(\rho)-Q_zg(\rho)+\Lambda(x,z),g(\rho)\bigr),
 \qquad \rho=y+Q(x,z),
\end{equation}
where $\Lambda$ is a homogeneous linear polynomial and $f(0)=g(0)=0$. If $Q_2$ is the quadratic homogeneous part of $Q$, expansion of the principal minors gives
\begin{align}\label{eq:quadratic-relation-derivative}
 \sigma_2(JH)
 &=f'(\rho)\bigl(Q_{xx}f(\rho)+Q_{xz}g(\rho)-\Lambda_x\bigr)\notag\\
 &\quad+g'(\rho)\bigl(Q_{xz}f(\rho)+Q_{zz}g(\rho)-\Lambda_z\bigr)\\
 &=\left.\frac{\dd}{\dd t}
 \bigl(Q_2(f(t),g(t))-\Lambda(f(t),g(t))\bigr)\right|_{t=\rho}.\notag
\end{align}
Substitution $t\mapsto \rho$ is injective. The polynomial being differentiated is constant and vanishes at zero. Thus
\begin{equation}\label{eq:quadratic-relation}
 Q_2(f(t),g(t))=\Lambda(f(t),g(t)).
\end{equation}

\smallskip\noindent\emph{Constant linear normalization.}
The shear $(x,y,z)\mapsto(x,y+Dx+Ez,z)$ removes the linear part of $Q$ from~\eqref{eq:quadratic-form}; the corresponding target change cancels the terms $-Df-Eg$. We may therefore assume $Q=Q_2$. If $Q_2=0$, relation~\eqref{eq:quadratic-relation} and independence of $f,g$ force $\Lambda=0$. The second component would then vanish, contradicting independence.

The binary quadratic form $Q_2$ cannot have rank two. Over $\overline K$, such a form is a product of independent linear forms. After evaluation at $(f,g)$, relation~\eqref{eq:quadratic-relation} would give
\[
  \bar{F}G=\alpha \bar{F}+\beta G,\qquad \bar{F}(0)=G(0)=0,
\]
with $\bar{F},G$ linearly independent over $ \overline K$. But
$(\bar{F}-\beta)(G-\alpha)=\alpha\beta$. If $\alpha\beta\neq0$, both factors are units of $\overline K[t]$. If $\alpha\beta=0$, one factor is zero. In either case, $\bar{F}$ or $G$ is constant, contradicting independence.

Thus $Q_2$ has rank one and can be written over $K$ as $\kappa L(x,z)^2$, with $\kappa\in K^*$ and $L$ a nonzero linear form. After a simultaneous linear change of $x,z$, write
\[
 Q_2=\kappa x^2,\qquad \Lambda=\alpha x+\beta z.
\]
Then $\kappa f^2=\alpha f+\beta g$. Necessarily $\beta\neq0$, and~\eqref{eq:quadratic-form} becomes
\begin{equation}\label{eq:precanonical}
 H=\left(f(\rho),\ \beta z-2\kappa xf(\rho)+\alpha x,
 \ \frac{\kappa}{\beta}f(\rho)^2-\frac{\alpha}{\beta}f(\rho)\right),
 \qquad \rho=y+\kappa x^2.
\end{equation}
The shear $z'=z+(\alpha/\beta)x$ removes the terms involving $\alpha$. Finally set
\[
 x'=x,\qquad y'=y/\kappa,\qquad z''=(\beta/\kappa)z'.
\]
With $P(t)=f(\kappa t)$, the transformed map is $\Ncal_P(x',y',z'')$. All these changes are defined over $K$, and $P(0)=0$.
\end{proof}

\end{document}
